\documentclass{article}
\usepackage[UTF8]{ctex}
\usepackage{amsmath,amssymb,amsfonts}
\usepackage{enumitem}
\usepackage{geometry}
\geometry{a4paper, margin=1in}
% === 错题本排版优化（不参与题目标识） ===
\setlength{\parindent}{0pt}
\setlength{\parskip}{0.35em plus 0.10em minus 0.08em}
\setlength{\emergencystretch}{3em}
\setlist{itemsep=0.25em, topsep=0.40em, parsep=0pt}
\allowdisplaybreaks[2]
% === 错题本排版优化结束 ===
\begin{document}
\section*{西北工业大学2024年高等代数真题}
\section*{171．西北工业大学2024年高等代数试卷}

\subsection{题目1}
一。设 \(A\) 是三阶复方阵，且满足关于迹的等式 \(\operatorname{tr}\left(A^{i}\right)=i(i=1,2,3)\) ，求行列式 \(\operatorname{det} A\) 。

\begin{itemize}
  \item[二．] 设 \(A\) 是 \(m \times n\) 矩阵，且 \(r(A)=r, B\) 是 \(m \times p\) 矩阵。
  \item[（1）] 给出矩阵方程 \(A X=B\) 有解的充要条件。
  \item[（2）] 若 \(X_{0}\) 是 \(A X=B\) 的一个特解，试求出 \(A X=B\) 的通解．
\end{itemize}

\subsection{题目2}
三。设 \(A, B\) 为数域 \(F\) 上的 \(m \times n, n \times s\) 矩阵，记 \(W=\left\{B X \mid X \in F^{s}, A B X=0\right\}\) 是 \(n\) 元（列）空间\(F^{n}\) 的子空间，证明： \(\operatorname{dim} W=r(B)-r(A B)\) ．

\subsection{题目3}
四．求满足 \(A^{*}=A\) 的所有 \(n(n \geq 2)\) 阶复方阵 \(A\) ，其中 \(A^{*}\) 是 \(A\) 的伴随矩阵。

\subsection{题目4}
五．设 \(J=\left(\begin{array}{rrrr}0 & 1 & & \\ & 0 & \ddots & \\ & & & 1 \\ & & & 0\end{array}\right)\) 是 \(n(n \geq 2)\) 阶 Jordan 块，求 \(J^{2}\) 的 Jordan 标准形．

\subsection{题目5}
六．求实二次型

\[
f\left(x_{1}, x_{2}, \cdots, x_{n}\right)=\left(x_{1}-\bar{x}\right)^{2}+\left(x_{2}-\bar{x}\right)^{2}+\cdots+\left(x_{n}-\bar{x}\right)^{2}
\]

的规范标准形，其中 \(\bar{x}=\frac{1}{n}\left(x_{1}+x_{2}+\cdots+x_{n}\right)\) ．

\subsection{题目6}
七．设 \(\eta\) 是 \(n\) 维 Euclid 空间 \(V\) 的一个单位向量．

\begin{itemize}
  \item[（1）] 证明： \(\mathcal{A}: \alpha \mapsto \alpha-2(\eta, \alpha) \eta(\forall \alpha \in V)\) 是 \(V\) 上的正交变换（也称镜面反射）．
  \item[（2）] 若 \(V\) 的正交变换 \(\mathcal{B}\) 有特征值1，且特征子空间 \(V_{1}\) 是 \(n-1\) 维的，证明： \(\mathcal{B}\) 是镜面反射。
\end{itemize}

\subsection{题目7}
八．在 \(F^{n \times n}\) 中，记

\[
W_{1}=\left\{A \in F^{n \times n} \mid \operatorname{tr}(A)=0\right\}, W_{2}=\left\{\operatorname{diag}\left\{x_{1}, x_{2}, \cdots, x_{n}\right\} \mid x_{1}=x_{2}=\cdots=x_{n} \in F\right\} .
\]

\begin{itemize}
  \item[（1）] 证明：\(W_{1}, W_{2}\) 为 \(F^{n \times n}\) 的子空间。
  \item[（2）] 求 \(W_{1}, W_{2}\) 的一个基和维数。
  \item[（3）] 证明：\(F^{n \times n}=W_{1} \oplus W_{2}\) ．
\end{itemize}

\subsection{题目8}
九。设线性空间 \(V\) 上的线性变换 \(\mathcal{A}\) 在一组基 \(\alpha_{1}, \alpha_{2}, \alpha_{3}\) 下满足

\[
\mathcal{A}\left(\alpha_{1}, \alpha_{2}, \alpha_{3}\right)=\left(\alpha_{1}, \alpha_{2}, \alpha_{3}\right)\left(\begin{array}{ccc}
3 & 4 & -5 \\
6 & 3 & -4 \\
6 & 4 & -6
\end{array}\right) .
\]

\begin{itemize}
  \item[（1）] 求 \(\mathcal{A}\) 的所有特征值和特征向量．
  \item[（2）] 求 \(\mathcal{B}=\mathcal{A}^{3}-5 \mathcal{A}\) 的一个非平凡不变子空间．
\end{itemize}

\subsection{题目9}
十．设 \(A, C\) 为 \(n\) 阶正定实对称矩阵，证明：\(A X+X A=C\) 存在唯一解 \(B\) ，且 \(B\) 也是正定矩阵。
\end{document}
